%% kapa: exam=1 round=34 year=2023 period=1 form=A subject=civil_law subject_ko=민법 seq=8 number=8
%% source: Q-Net 감정평가사 기출문제 (cst00309) artlSeq=5211789 file=raw/1008/5211789/2023년도 제34회 감정평가사 1차시험 1교시.pdf
%% conversion: kordoc 4.24.51 -> tools/gen1.py (auto); review=proofread-sonnet; flags=none
\begin{question}{8}
물건에 관한 설명으로 옳지 않은 것은? (다툼이 있으면 판례에 따름)
\begin{choices}
  \item 주물의 구성부분은 종물이 될 수 없다.
  \item 1필의 토지의 일부는 분필절차를 거치지 않는 한 용익물권의 객체가 될 수 없다.
  \item 국립공원의 입장료는 법정과실이 아니다.
  \item 주물과 장소적 밀접성이 인정되더라도 주물 그 자체의 효용과 직접 관계가 없는 물건은 종물이 아니다.
  \item 저당권 설정행위에 “저당권의 효력이 종물에 미치지 않는다”는 약정이 있는 경우, 이를 등기하지 않으면 그 약정으로써 제3자에게 대항할 수 없다.
\end{choices}
\end{question}
\answer{2}
